% Generated draft: requires mathematical and source review.
\documentclass[11pt]{article}
\usepackage[margin=25mm]{geometry}
\usepackage{fontspec}
\setmainfont{Latin Modern Roman}
\usepackage{amsmath,amssymb,mathtools,unicode-math}
\setmathfont{Latin Modern Math}
\usepackage{enumitem,needspace,xurl,hyperref}
\setlength{\parindent}{0pt}
\setlength{\parskip}{5pt}
\setlength{\emergencystretch}{3em}
\newcommand{\N}{\mathbb{N}}
\newcommand{\Z}{\mathbb{Z}}
\newcommand{\Q}{\mathbb{Q}}
\newcommand{\R}{\mathbb{R}}
\newcommand{\C}{\mathbb{C}}
\newcommand{\F}{\mathbb{F}}
\newcommand{\A}{\mathbb{A}}
\newcommand{\PP}{\mathbb{P}}
\newcommand{\E}{\mathbb{E}}
\newcommand{\eps}{\varepsilon}
\newcommand{\dd}{\,\mathrm{d}}
\begin{document}
\section*{UnsolvedMath ID: 30004752}
\textbf{Source problem code:} catalogue entry
\subsection{Definitions and mathematical statement}
Let \(E/\Q\) be an elliptic curve without complex multiplication, meaning \(\operatorname{End}_{\overline{\Q}}(E)=\Z\). For each prime \(\ell\), let
\[
T_\ell(E)=\varprojlim_n E[\ell^n](\overline{\Q}),\qquad
\rho_{E,\ell^\infty}:\operatorname{Gal}(\overline{\Q}/\Q)\longrightarrow\operatorname{Aut}_{\Z_\ell}(T_\ell(E))\cong\operatorname{GL}_2(\Z_\ell)
\]
be the \(\ell\)-adic Galois representation. Its image is considered up to conjugacy, since the displayed identification depends on a choice of basis.

\textbf{Intended completeness question.} Is the source's known list of \(22\) exceptional \(\ell\)-adic Galois images complete for non-CM elliptic curves over \(\Q\)? More explicitly, the intended quantification is over every non-CM elliptic curve \(E/\Q\) and every prime \(\ell\): whenever the image of \(\rho_{E,\ell^\infty}\) is exceptional in the source's sense, must it be one of the listed images, with the prime and conjugacy conventions of that list? The supplied record does not enumerate these \(22\) images or define the criterion for being exceptional, so this is a qualified formulation rather than a fully specified classification statement.

\textbf{Supplied counterexample reduction.} The record additionally asks whether every counterexample to completeness must arise from an unexpected rational point on a curve in the following collection:
\[
\bigl\{X_{\mathrm{ns}}^+(N):N\in\{33,5^2,7^2,11^2\}\bigr\}
\;\cup\;
\bigl\{X_{\mathrm{ns}}^+(\ell):\ell\text{ prime},\ \ell\geq17\bigr\}
\;\cup\;\{Y_1,Y_2\}.
\]
Here \(X_{\mathrm{ns}}^+(N)\) is the source's notation for the relevant nonsplit Cartan normalizer modular curve, and \(Y_1,Y_2\) are placeholders for the source's two level-\(49\) modular curves of genus \(9\). Their defining subgroups, models, and maps are not supplied. The terms ``unexpected'' and ``arise from'' are also unspecified, including which rational points are excluded or already known and how the remaining points produce exceptional images. The level \(33\) is retained exactly as written in the clean statement.

The record calls the counterexample reduction equivalent to completeness. A restriction on where counterexamples can arise does not itself exclude counterexamples; an additional assertion ruling out the relevant unexpected points or their contribution would be needed to establish that equivalence. The original statement asserts completeness, whereas the clean statement poses it as a question. This draft retains the question and records the unresolved logical qualification.
\subsection{Short English statement}
For each non-CM elliptic curve over \(\Q\) and each prime \(\ell\), ask whether every exceptional \(\ell\)-adic Galois image belongs to the source's list of \(22\). The supplied record relates possible counterexamples to rational points on specified modular curves, but omits the list, the meaning of exceptional and unexpected, and the identities of two curves. Its claimed equivalence requires clarification.
\subsection{Sources}
\begin{itemize}
\item Local UnsolvedMath registry: supplied record with permanent ID \(30004752\), title ``Exceptional Galois Images of Non-CM Elliptic Curves,'' source problem code \texttt{catalogue entry}, in the Oberwolfach Reports Open Problems collection. No individual proposer is supplied; the proposed year is recorded as \(2021\).
\item Supplied report citation: \emph{Explicit Methods in Number Theory (2022)}, Oberwolfach Reports, DOI \texttt{10.4171/owr/2021/34}, pp.\ 1805--1858. \url{https://doi.org/10.4171/owr/2021/34}. The record describes these reports as documenting mathematical workshops and open-problem discussions. The citation's year \(2022\) and the record's proposed year \(2021\) are preserved without reconciliation.
\item Supplied literature reference, labeled ``Rouse--Sutherland--Zureick-Brown'': \url{https://arxiv.org/abs/2106.11141}.
\item Supplied literature reference, labeled ``Published classification'': \url{https://www.cambridge.org/core/services/aop-cambridge-core/content/view/D5BC92F9949B387570A7D764635B6AC8/S205050942200038Xa.pdf/dollarell-dollar-adic-images-of-galois-for-elliptic-curves-over-dollarmathbb-qdollar-and-an-appendix-with-john-voight.pdf}.
\item Literature assessment dated exactly \texttt{2026-08-22}: the supplied curation labels the problem open and states that completeness of the exceptional-image list and remaining rational-point questions on several nonsplit Cartan modular curves and related levels are unresolved. This is supplied dated curation derived from the OWR material, not an independently verified literature assessment.
\item The registry labels the statement \texttt{corrected\_verified} and says it was checked against the official Oberwolfach report. This is supplied verification metadata; no independent source visit or mathematical verification is asserted here.
\end{itemize}
\end{document}
